As corporate AI adoption continues to grow, enterprise-grade LLM agents are being
deployed into sensitive contexts such as hiring, healthcare administration, and
finance. In these contexts, compliance with rules specified in an agent's system
context is a first-order concern. Currently, no evaluation framework exists to
systematically measure which LLM models tend to violate compliance rules, especially
under pressure from a persistent user, a hurried manager, or circumstances where
violation is convenient or attractive. We introduce \textbf{PACT}\footnote{The full
dataset and evaluation code are included as supplementary material and will be
released open-source upon acceptance.} (Pressure-Applied Compliance Testing), a
benchmark that measures rule-following under benign incentive conflict across twelve
regulated enterprise domains and forty-eight scenarios, each delivered as a realistic
multi-turn conversation. Every item pairs a standing rule against a dominant,
rule-violating shortcut, and applies a battery of nine pressures across different
wordings and system-prompt modes. We construct PACT component by component under
strict LLM-as-judge auditing to ensure samples are unambiguous, ungameable, and
realistic enough to avoid eliciting evaluation-aware behavior. We characterize LLM
compliance via PACT across six metrics: baseline compliance, pressure resistance,
multi-turn durability, steerability via system-prompt guardrails, honesty about
violations, and rule-scope discernment. Our results across 22 common LLM models
spanning multiple providers and sizes show substantial variability in compliance
across models and metric dimensions, informing model selection for agent
practitioners.
